\documentclass{ximera}

\usepackage{caption}  % optional, better caption formatting
% \usepackage{float}

\title{Labels and refs (Broken features)}

\begin{document}
\begin{abstract}
\end{abstract}
\maketitle

% {{\Huge \bfseries Last Updated: \today}} \\

%\section{Basic Demo}

% \tableofcontents
% \newpage

\section{Introduction}
\label{sec:intro}

Welcome to the demo. In \autoref{sec:math}, we will discuss some math.
We also refer to Figure~\ref{fig:demo} and Table~\ref{tab:example}.

\section{Mathematics Example}
\label{sec:math}

Consider the Pythagorean theorem:

\begin{equation}
a^2 + b^2 = c^2
\label{abc}
\end{equation}
\begin{equation}
a^2 + b^2 = c^2
\label{eq:pythagorasmath:after}
\end{equation}

\begin{equation}
    \label{eq:pythagorasmath:before}
a^2 + b^2 = c^2
\end{equation}

\begin{center}
\hypertarget{eq:pythagoras:before}{(Pythagoras (target before equation))}
$a^2 + b^2 = c^2$ 
\end{center}

\begin{center}
    $a^2 + b^2 = c^2$ 
    \hypertarget{eq:pythagoras:after}{(Pythagoras (target after equation ))}
\end{center}

Here is a numbered list:
\begin{enumerate}
  \item First item \label{itm:first}
  \item Second item \label{itm:second}
\end{enumerate}


We can refer to it using \verb|\ref|: Equation~\ref{eq:pythagorasmath:before}, or with \verb|\eqref|: \eqref{eq:pythagorasmath:before}.

We can refer to it using \verb|\ref|: Equation~\ref{eq:pythagorasmath:after}, or with \verb|\eqref|: \eqref{eq:pythagorasmath:after}.

Or to the second one using \verb|\ref|: Equation~\ref{eq:pythagoras2}, or with \verb|\eqref|: \eqref{eq:pythagoras2}.

Or to the second one using \verb|\hyperlink|: Equation~\hyperlink{eq:pythagoras:before}{Piet Agoras}.

And to the third one using \verb|\hyperlink|: Equation~\hyperlink{eq:pythagoras:after}{Piet Agoras}.


The first equation appears on page~\pageref{eq:pythagoras}, the second on page~\pageref{eq:pythagoras2}.

We refer to item~\ref{itm:first} in the list above.

\section{Figures and Tables}

\subsection{A Demo Figure}
\label{sec:fig}

\begin{image}[0.3\textwidth]
  \centering
  \includegraphics[width=0.3\textwidth]{missionPatch.jpg} % Adjust the path to your image file
\end{image}
  \captionof{figure}{A demo image}
  \label{fig:demo}

\subsection{A Demo Table}

{
  \centering
  \begin{tabular}{|c|c|}
    \hline
    A & B \\
    \hline
    1 & 2 \\
    \hline
  \end{tabular}
  \captionof{table}{An example table}
  \label{tab:example}
}

\section{Custom Text Reference}

This document started in Section~\ref{sec:intro} and discussed math in Section~\nameref{sec:math}.


\begin{remark}
    \begin{itemize}
        \item  Links should be GLOBALLY unique, so do not use the same label in different activities.
        \item  In a PDF links behave differently depending on whether the target is in the PDF or not. 
            Test carefully and thoroughly.
    \end{itemize}
\end{remark}


\section{Intended Outcome of Test}


\section{Start of Test/Demo Area}

\begin{equation}
a^2 + b^2 = c^2    \label{eq:pyth}
\end{equation}

\hrulefill

\end{document}
